% Aannames
% Not a real \chapter (it's a hand-styled heading like Register's), so it
% needs its own Inhoud/PDF-bookmark entry, same as imakeidx's intoc=true
% does for Register.
\phantomsection\addcontentsline{toc}{chapter}{Aannames}
\vspace*{0.4cm}
{\handfont\fontsize{60}{60}\selectfont\color{terracotta}Aannames\par}
\vspace{8pt}{\color{terracotta}\rule{2.6cm}{0.5pt}\;\raisebox{-0.2ex}{$\scriptstyle\blacklozenge$}}\par
\vspace{1.4em}

\lettrine{E}{en} kookboek zit vol keuzes die nergens met zoveel woorden
worden uitgesproken: welke ovenstand hoort erbij, en wat is er precies
bedoeld als er gewoon ``boter'' staat? Hieronder staan de aannames die we
in dit boek overal aanhouden, zodat een recept ook duidelijk is op de
punten waar de tekst zelf niets over zegt.

\blockrule{Oven}
Staat er bij een recept ``Verwarm de oven voor'', dan is dat steeds
boven- en onderwarmte, tenzij het recept nadrukkelijk een andere stand
noemt. Zet de oven op tijd aan: de meeste recepten gaan pas de oven in als
hij al op temperatuur is.

\blockrule{Basisingrediënten}
\begin{ingredients}
\ingb{Eieren zijn maat M, tenzij een recept iets anders noemt}
\ingb{Roomboter of boter is ongezouten, tenzij een recept nadrukkelijk gezouten boter of margarine vraagt}
\ingb{Peper is versgemalen zwarte peper, in plaats van kant-en-klaar gemalen}
\end{ingredients}
Staat er bij een recept los ``zout en peper, naar smaak'' zonder vaste
hoeveelheid, dan is dat een kwestie van proeven, niet van een vergeten
hoeveelheid.

\vspace{0.6em}
{\footnotesize\itshape\color{inkfaint}Dit zijn de aannames die we zelf
aanhouden, geen harde regels. Noemt een recept iets nadrukkelijk anders,
volg dan gewoon het recept.\par}
